% from ega0-23-fr.tex
\oldpage[0\textsubscript{IV-1}]{216}

que $B$ est un anneau japonais, autrement dit, on peut se borner au cas où
$A$ est en outre \emph{régulier}. Raisonnons par récurrence sur
$n=\dim(A)$, le théorème étant évident pour $n=0$. Supposons donc $n>0$, et,
en désignant par $\mathfrak m$ l'idéal maximal de $A$, soit $x$ un élément de
$\mathfrak m-\mathfrak m^2$~; on sait (17.1.8) que $A/xA$ est un anneau
régulier, et en outre ((17.1.7) et (16.3.4)) que
$\dim(A/xA)=n-1$~; de plus, l'idéal $xA$ est fermé dans $A$
(0\textsubscript{I}, 7.3.5), donc $A/xA$ est complet. En vertu de l'hypothèse
de récurrence, $A/xA$ est donc un anneau japonais, et il résulte alors de
(23.1.3) qu'il en est de même de $A$.
\end{proof}

\begin{corollary}[23.1.6]
\label{0.23.1.6-fr}
Soient $A$ un anneau local noethérien intègre complet, $K$ son corps des
fractions, $K'$ une extension finie de $K$~; alors la fermeture intégrale
$A'$ de $A$ dans $K'$ est un anneau local complet.
\end{corollary}

\begin{proof}
On sait déjà par (23.1.5) que $A'$ est une $A$-algèbre finie, donc un anneau
semi-local noethérien complet, et par suite un \emph{composé direct}
d'anneaux locaux~; mais comme $A'$ est intègre, c'est nécessairement un
anneau local.
\end{proof}

\begin{proposition}[23.1.7]
\label{0.23.1.7-fr}
Soient $A$ un anneau local noethérien intègre, $K$ son corps des fractions,
$K'$ une extension finie de $K$, $A'$ la fermeture intégrale de $A$ dans
$K'$. On suppose que le complété $\widehat A$ est réduit, et l'on désigne par
$R$ son anneau total des fractions.
\begin{enumerate}
  \item[(i)] Si l'anneau $R\otimes_KK'$ est réduit (ce qui aura lieu en
  particulier lorsque $K'$ est une extension séparable de $K$, $R$ étant
  composé direct de corps, extensions de $K$ (Bourbaki, \emph{Alg.},
  chap.~VIII, §~7, n\textsuperscript{o}~3, cor.~1 du th.~1)), alors $A'$ est
  un $A$-module de type fini.
  \item[(ii)] Si en particulier $R$ est une $K$-algèbre séparable, alors $A$
  est un anneau japonais.
\end{enumerate}
\end{proposition}

\begin{proof}
\emph{(i)} Posons pour simplifier $A_1=\widehat A$,
$A'_1=A'\otimes_AA_1$, $K'_1=K'\otimes_AA_1$. Comme $A_1$ est un
$A$-module plat (0\textsubscript{I}, 7.3.5), $A'_1$ s'identifie à un
sous-anneau de $K'_1$ et est évidemment entier sur $A_1$. D'autre part, si
l'on pose $K_1=K\otimes_AA_1$, on peut écrire
$K'_1=K'\otimes_KK_1$~; comme $A$ est intègre et que $A_1$ est un $A$-module
plat, tout élément $\neq0$ de $A$ est $A_1$-régulier
(0\textsubscript{I}, 6.3.4)~; comme $K=S^{-1}A$, où $S=A-\{0\}$, on a
$K_1=S^{-1}A_1$, et la remarque précédente prouve que $K_1$ s'identifie à un
sous-anneau de l'anneau total des fractions $R$ de $A_1$. Puisque tout
$K$-module est plat, $K'_1$ s'identifie à un sous-anneau de $K'\otimes_KR$,
qui par hypothèse est réduit et est une $R$-algèbre finie. Désignons par
$\mathfrak q_i$ ($1\leqslant i\leqslant n$) les idéaux premiers minimaux de
$A_1$, par $B_i$ l'anneau intègre $A_1/\mathfrak q_i$, par $L_i$ le corps des
fractions de $B_i$, de sorte que $R$ est composé direct des $L_i$, et
$K'\otimes_KR$ composé direct des $K'\otimes_KL_i$~; ces dernières
$K$-algèbres, étant réduites, sont des composées directes de corps,
extensions finies des $L_i$. Comme $B_i$ est un anneau local intègre
\emph{complet}, le théorème de Nagata (23.1.5) prouve que la fermeture
intégrale de $B_i$ dans $K'\otimes_KL_i$ est un $B_i$-module de type fini, et
\emph{a fortiori} un $A_1$-module de type fini~; comme tout élément de
$K'\otimes_KL_i$ qui est entier sur $A_1$ est aussi entier sur $B_i$ (étant
annulé par $\mathfrak q_i$), on en conclut finalement que la fermeture
intégrale de $A_1$ dans $K'\otimes_KR$ est un $A_1$-module de type fini. Mais
comme $A'_1$ est contenu dans cette fermeture intégrale et que $A_1$ est
noethérien, $A'_1$ est aussi un $A_1$-module de type fini. Enfin, comme $A_1$
est un $A$-module fidèlement plat (0\textsubscript{I}, 7.3.5), on en conclut
que $A'$ est un $A$-module de type fini (Bourbaki, \emph{Alg. comm.},
chap.~I, §~3, n\textsuperscript{o}~6, prop.~11).

